% ============================================================
% ADVISORY OUTLINE — DRAFT ONLY
% ============================================================
% Purpose: Committee check-in document (target: May 2026).
% Shows the committee how theory, making, and ST entangle
% across the Composer's Note (introduction) and all three
% Movements. Demonstrates command of the full score even
% while Movement II receives deepest development.
%
% This file is DRAFT ONLY. It is guarded by \ifdraft in
% main.tex and will NOT appear in the final dissertation.
% ============================================================

% ------------------------------------------------------------
% HOW TO READ THIS OUTLINE
% ------------------------------------------------------------

\section*{How to Read This Outline}

This outline maps the intellectual architecture of the dissertation as it stands in April 2026, written toward a September defense. It is organized to show you three things at once: (1) what theoretical and literary sources anchor each section; (2) how the making practice entangles with those sources at specific moments rather than in a prior chapter; and (3) where \textit{Speculocultural Technopoiesis} (ST) is being applied, tested, and extended across each movement.

Theory arrives when the material demands it. All knowledge is situated. This outline tries to make that entanglement visible and followable.

\textbf{Scope note (April 2026):} Movement II (\textit{Did Sun Ra Fear Cryosleep?} / water) is the primary site of deep analysis, conference presentation, and collaborative workshop documentation. Movements I and III are developed as adjacent works: architecturally complete, theoretically grounded, and understood as a continuation of practice extending beyond the dissertation. All three are present in the score, though they carry different analytical weight.

\medskip
\noindent\textbf{ST Domain Key:}
\begin{itemize}[noitemsep]
    \item \textbf{FS} --- Fugitive Speculation
    \item \textbf{MM} --- Material Memory
    \item \textbf{SW} --- Sensory Worlding
    \item \textbf{IW} --- Intimate Witnessing
\end{itemize}

\noindent\textbf{ST Cycle:} Unsettling $\to$ Remixing $\to$ Embodying $\to$ Transmitting

\bigskip

% ------------------------------------------------------------
% PART I: COMPOSER'S NOTE (Introduction)
% ------------------------------------------------------------

\section*{Part I: Composer's Note}
\addcontentsline{toc}{section}{Part I: Composer's Note (Introduction)}

The Composer's Note serves as the dissertation's introduction. It establishes positionality, articulates ST as a framework, and explains the organizing metaphors of the score, but in a register consistent with the form. You are addressed as performers and listeners, not as evaluators. The \textit{Note on Structure} chapter, which precedes the Composer's Note in the full text, makes the institutional argument for that form explicitly, so the Composer's Note can proceed without interruption.

\medskip

\subsection*{The Drone}

\noindent\textbf{Function:} Opens the dissertation in sound rather than argument. Sound arrives before any instrument shapes it. This section introduces the perceptual constant, the drone, that runs beneath all three movements.

\noindent\textbf{Key Literature:}
\begin{itemize}[noitemsep]
    \item Grimshaw \& Garner, \textit{Sonic Virtuality}: sound as a field of potential, realized through the conditions of its mediation. Transduction as a culturally shaped conversion of sonic energy.
\end{itemize}

\noindent\textbf{Theory--Making Entanglement:} The drone is the perceptual substructure of all three instruments. Sonic virtuality gives this a name: the field of potential that exists before any one instrument conditions it into a specific sound. ST begins at the moment a practitioner becomes conscious of what was already there, already shaping what could be heard. Manipulated and formed by the defaults of the system.

\noindent\textbf{ST Connection:} The drone is the pre-unsettled condition.

\bigskip

\subsection*{Who Is Making This}

\noindent\textbf{Function:} Positionality. ST is a \textit{Black Sonic} Speculocultural Technopoiesis. The cultural grounding is the research site, and one that does not precede the argument. The "specific-to-universal" structure: beginning in a particular place, deeply enough, generates principles that travel.

\noindent\textbf{Key Literature:}
\begin{itemize}[noitemsep]
    \item Wynter, ``The Ceremony Found'': the autopoetic act of re-envisioning the human beyond Western epistemic closure; ST as a sonic ceremony-finding \parencite{Wynter_2015}.
\end{itemize}

\noindent\textbf{Theory--Making Entanglement:} My body is the primary instrument here: my formation as a musician, my fluencies (both technological and material) my public practice, my dual position as artist and scholar. Wynter helps to put a name to the face of what that dual position enacts. The ceremony finds itself through the making.

\bigskip

\subsection*{A Trombone in Third Position}

\noindent\textbf{Function:} Origin story and core epistemological claim. My assistant band director once explained to me that my Yamaha trombone plays slightly sharp in third position compared to a Bach or a Shires. Each maker designed the instrument around a distinct philosophy of sound. and to another degree to the tools and worldviews of those tools that they in turn used to craft their own. A lineage of worldviews and learned constraints. That detail was a lesson in epistemology I didn't have words for yet: instruments carry the intentions of their makers. To modify a tool is to engage in an epistemic act that is conversation or argument with its original maker.

\noindent\textbf{Key Literature:}
\begin{itemize}[noitemsep]
    \item Grimshaw \& Garner, \textit{Sonic Virtuality}: transduction and refraction as the mechanism of tool-mediated perception.
    \item Eshun, \textit{More Brilliant Than the Sun}: sonic fiction as a speculative apparatus; temporal reconfiguration through sound \parencite{Eshun_1999}.
\end{itemize}

\noindent\textbf{Theory--Making Entanglement:} The story arrives before the theory. The theory helps to articulate something the story already knew. That sequencing critique-practice-theory-critique (repeating), is the research-creation method in miniature.

\bigskip

\subsection*{Making the Score: On the Structure of This Dissertation}

\noindent\textbf{Function:} Articulates the ST cycle (Unsettling, Remixing, Embodying, Transmitting) and the four cross-cutting domains (FS, MM, SW, IW). Names the two simultaneous temporalities running through all three works: elemental/geological (earth $\to$ water $\to$ body) and historical/speculative (buried past $\to$ suspended passage $\to$ present-tense politics). Introduces Julius Eastman's \textit{Femenine}, and his Organic Music, as the structural model: drone as perceptual constant, ST cycle as the developmental cell.

\noindent\textbf{Key Literature:}
\begin{itemize}[noitemsep]
    \item Benjamin, \textit{Race After Technology}: ``the road to inequity is paved with technical fixes'' \parencite{Benjamin_2019a}; surface-level interventions leave the epistemic foundation untouched.
    \item Wynter, ``Unsettling the Coloniality of Being'': the over-representation of ``Man'' as universal through supposedly neutral logics \parencite{Wynter_2003}.
    \item Manning \& Massumi, \textit{Thought in the Act}: thinking-feeling; the inseparability of conceptual and bodily engagement in making \parencite{Manning_Massumi_2014}.
    \item Loveless, \textit{How to Make Art at the End of the World}: research-creation as inquiry driven by practice rather than predetermined frameworks \parencite{Loveless_2019}.
    \item Eastman, \textit{Femenine}: structural model for how a drone and a developmental cell sustain a work across duration without resolving it.
\end{itemize}

\noindent\textbf{Theory--Making Entanglement:} By the time you reach the formal description of ST here, you have already seen it enacted: in the trombone story, in the positionality, in the drone. Each element of the cycle and each domain arrives named, with the practice already behind it. The framework should read as description, not prescription; as something that clarifies what was already happening.

\bigskip

\subsection*{How to Hear This}

\noindent\textbf{Function:} Listening guide. Presents the three movements and their tempos: geological (Ground), chemical (Hold), human (Touch). Explains that theory arrives in the material's service throughout, as it is being summoned by the making, committed and accountable to it. It describes the arc across all three movements: subliminal (felt, Ground) $\to$ emerging (witnessed, Hold) $\to$ explicit (participated, Touch). Prepares you for what the score does and why it does it that way.

\noindent\textbf{Key Literature:}
\begin{itemize}[noitemsep]
    \item Ihde, \textit{Technology and the Lifeworld}: the technology-world relations; named here as the condition all three instruments share \parencite{Ihde_1990}.
\end{itemize}

\noindent\textbf{Theory--Making Entanglement:} By the time this section arrives, I already know what each movement will find. Not because it has been predetermined, but because fo the anticipated directions grounded in two-plus years of making. Framed through the enabling constraints of this work. Ihde's hermeneutic relation is the thread: every instrument in this score shapes what can be heard through its design, and the research is built on making that shaping conscious.

\bigskip

% ------------------------------------------------------------
% PART II: MOVEMENT I — GROUND
% ------------------------------------------------------------

\section*{Part II: Movement I --- Ground/Earth (\textit{Buffalo Soldier}, earth)}
\addcontentsline{toc}{section}{Part II: Movement I --- Ground}

\noindent\textbf{Tempo:} Geological. The pace of sedimentation, burial, and what the earth holds without announcing it.\\
\noindent\textbf{Form:} Excavation log. Layers that go deeper rather than forward.\\
\noindent\textbf{Material:} Wooden substrate boxes, earth, rock, soil, piezoelectric sensors, electret microphones, Max/MSP granular processing chain.\\
\noindent\textbf{ST Domain Foregrounded:} Material Memory (MM).\\
\noindent\textbf{Motif Recognition:} Subliminal. Felt in the body, not consciously heard.\\
\noindent\textbf{Listener Relationship:} Felt.\\
\noindent\textbf{Development Status:} Adjacent work. Architecturally complete, theoretically grounded. Prose to be developed from existing scaffolding.

\medskip

\subsection*{Surface: The Desire and the Unsettling}

\noindent\textbf{Literature/Theory:}
\begin{itemize}[noitemsep]
    \item Estes, \textit{Our History Is the Future}: the deliberate near-extinction of the bison as colonial weapon; the triangulated histories encoded in the name ``Buffalo Soldier'' — Indigenous naming, Black body, colonial violence \parencite{Estes_2019}.
    \item Benjamin, \textit{Race After Technology}: default reverb and spatialization algorithms assume neutral acoustic space. Space in these tools carries no memory, no weight, no history.
\end{itemize}

\noindent\textbf{Making Entanglement:} The work begins in my family's home, from my father's perspective: paintings of Buffalo Soldiers, a Colt revolver replica, cavalry swords. Objects that carry doubled history. The unsettling happens the first time standard Max/MSP reverb produces a sonic space with no memory of what the substrate materials hold. That gap between what the earth carries and what the tool can process is the research problem, encountered in practice.

\noindent\textbf{ST Phase:} Unsettling.

\medskip

\subsection*{First Layer: The Ground as Archive}

\noindent\textbf{Literature/Theory:}
\begin{itemize}[noitemsep]
    \item McKittrick, \textit{Demonic Grounds}: Black geographies; land as a site of contested demarcation, extraction, and burial \parencite{McKittrick_2006}.
    \item Bennett, \textit{Vibrant Matter}: rocks and soil are not passive; thing-power; the substrate co-authors the sonic output \parencite{Bennett_2010}.
\end{itemize}

\noindent\textbf{Making Entanglement:} The wooden substrate box is the earth that holds memory and the container that confines bodies. It is coffin, cargo hold, reservation, plantation, and boundary; all simultaneous. The piezoelectric sensor listens to the earth it rests in with meticulous anticipation, waiting for the ground to confirm what it already believes it will find. Bennett's thing-power is operational here: smooth river rock, dense soil, and sand each produce distinct sonic characters, and the material makes decisions the patch did not anticipate.

\noindent\textbf{ST Phase:} Unsettling $\to$ Remixing. Identifying what the tool cannot hold, then beginning to redesign the processing chain around it.

\medskip

\subsection*{Second Layer: The Encounter That Has No Archive}

\noindent\textbf{Literature/Theory:}
\begin{itemize}[noitemsep]
    \item Hartman, ``Venus in Two Acts'' / Critical Fabulation: filling archival silence through speculative method; the archive was never built to hold this encounter \parencite{Hartman_2008}.
\end{itemize}

\noindent\textbf{Making Entanglement:} A Buffalo Soldier meets an American buffalo. They recognize each other. Both of them commodified, both of them nearly driven to extinction, both carrying the weight of the same genocidal logic. That meeting has no archive. The granular synthesis chain enacts the fabulation: sounds torn apart and reconstituted, tracking the violence of extraction and the slow accumulation of memory. Spatialization follows migration route data and forced march trajectories. The spatial logic is the argument.

\noindent\textbf{ST Phase:} Remixing. The modification is ideological before it is technical.

\medskip

\subsection*{Third Layer: The Right to Not Be Legible}

\noindent\textbf{Literature/Theory:}
\begin{itemize}[noitemsep]
    \item Glissant, \textit{Poetics of Relation}: opacity as a right; the refusal of full legibility to dominant systems of knowledge \parencite{Glissant_1997}.
    \item Ihde, \textit{Technology and the Lifeworld}: the hermeneutic relation applied to the piezo as a designed listener; adjusting the sensor's sensitivity thresholds is an epistemic act \parencite{Ihde_1990}.
\end{itemize}

\noindent\textbf{Making Entanglement:} The piezo picks up everything without hierarchy or understanding of boundary or history. The sonic ambiguity between those materials is the argument. Both the Buffalo Soldier and the buffalo were rendered legible only as functions of colonial utility. The instrument produces sound that resists clean categorization, and that resistance is deliberate.

\noindent\textbf{ST Phase:} Embodying. Testing whether the modification produces the experience of opacity that the cultural framework calls for.

\medskip

\subsection*{Fourth Layer: The Instrument Built from This}

\noindent\textbf{Literature/Theory:}
\begin{itemize}[noitemsep]
    \item Barad, \textit{Meeting the Universe Halfway}: the substrate box as apparatus; the earth-sound emerges through intra-action of materials, sensors, processing chain, and encoded cultural logic — not from any single element \parencite{Barad_2007}.
    \item Franinović \& Serafin, \textit{Sonic Interaction Design}: the PebbleBox as precedent; tactile interaction as primary sonic generation modality \parencite{Franinovic_Serafin_2013}.
\end{itemize}

\noindent\textbf{Making Entanglement:} Technical documentation: substrate box construction, piezo wiring and sensitivity calibration, electret microphone placement, the Max/MSP granular processing chain, multi-channel spatialization informed by migration data. The Somax2 corpus draws on field hollers, work songs, and bugle calls. Training the AI to find that material in earth sounds. The full ST cycle is documented here.

\noindent\textbf{ST Phase:} Transmitting. Documented procedures, parameter decisions, and the cultural reasoning behind them circulate with the work.

\medskip

\subsection*{Deep Layer: What Ground Hears}

\noindent\textbf{Making Entanglement:} Synthesis of the movement and what this instrument reveals. The leitmotif: \textit{When Johnny Comes Marching Home}, squashed and stretched to varying degrees of the representation of geological timescales by layer. Felt with/in the body before it is heard. The movement closes with the question that opens the next: \textit{If the ground holds this memory, if earth is archive, then what does water carry? What happens when the earth, the substrates of the body, enters the ocean?}

\noindent\textbf{ST Assessment:} What the cycle produced, what remains unresolved, what the next iteration would address.

\bigskip

% ------------------------------------------------------------
% PART III: MOVEMENT II — HOLD (PRIMARY)
% ------------------------------------------------------------

\section*{Part III: Movement II --- Hold (\textit{Did Sun Ra Fear Cryosleep?}, water)}
\addcontentsline{toc}{section}{Part III: Movement II --- Hold (PRIMARY)}

\adnote[title={Primary Movement}]{Movement II is the primary site of deep analysis, collaborative workshop documentation, and conference presentation (SfNC June 2, UCR Co-Creativity June 5--6, 2026). The outline below represents the full intellectual architecture of this movement. Sections are being written in order from April 2026 onward.}

\noindent\textbf{Tempo:} Chemical. The pace of dissolution, transformation in suspension.\\
\noindent\textbf{Form:} Ship's log / dissolution record. Entries that track transformation across time; sections that do not resolve.\\
\noindent\textbf{Material:} Water vessel, hydrophones, dissolving artifacts (salt, sugar, wax, clay/plaster), NASA sonification samples, Max/MSP spectral analysis and granular synthesis chain, conductivity sensor (planned).\\
\noindent\textbf{ST Domain Foregrounded:} Fugitive Speculation (FS).\\
\noindent\textbf{Motif Recognition:} Emerging. Recognizable in outline.\\
\noindent\textbf{Listener Relationship:} Witnessed.\\
\noindent\textbf{Development Status:} Primary movement. Deep writing, collaborative workshops (max 3 sessions, 1 collaborator confirmed), and conference deliverables all centered here.

\medskip

\subsection*{Entry: The Question}

\noindent\textbf{Literature/Theory:}
\begin{itemize}[noitemsep]
    \item Eshun, \textit{More Brilliant Than the Sun}: sonic fiction as a speculative apparatus; the speculative question as a research instrument that opens rather than closes \parencite{Eshun_1999}.
    \item Womack, \textit{Afrofuturism}: technology, liberation, and the construction of Black futurity; Sun Ra's mythology as Afrofuturist practice \parencite{Womack_2013}.
\end{itemize}

\noindent\textbf{Making Entanglement:} ``Did Sun Ra fear cryosleep?'' is the research instrument. The question re-inserts Black subjectivity — fear, willingness, choice — into a visual motif (\textit{Foundation}, \textit{Alien}, \textit{Interstellar}) that routinely evacuates it. The Middle Passage and the cryo-pod share an architectural logic: containment of Black bodies in fluid, in transit, across impossible space-time. That convergence is the design problem this work is built on.

\noindent\textbf{ST Phase:} Unsettling. Encountering the visual and sonic archive that treats Black bodies in suspension as backdrop.

\medskip

\subsection*{Log Entry 1: The Wake}

\noindent\textbf{Literature/Theory:}
\begin{itemize}[noitemsep]
    \item Sharpe, \textit{In the Wake: On Blackness and Being}: the ship, the hold, the wake, the weather; wake work as a methodology of tending to the dead while the living continue; the Atlantic as both grave and archive \parencite{Sharpe_2016}.
    \item Ihde, \textit{Technology and the Lifeworld}: the hydrophone as hermeneutic relation; what it makes legible is shaped by what it was designed to detect and by what the processing chain encodes \parencite{Ihde_1990}.
\end{itemize}

\noindent\textbf{Making Entanglement:} Submerging the hydrophone is the methodology. Listening in the wake means listening in the sonic residue of bodies that passed through water — and the default signal processing that ``cleans'' those recordings erases that residue. That erasure is the unsettling. The modification begins at the processing chain: amplifying roughness and instability rather than suppressing it.

\noindent\textbf{ST Phase:} Unsettling $\to$ Remixing.

\medskip

\subsection*{Log Entry 2: What the Water Holds}

\noindent\textbf{Literature/Theory:}
\begin{itemize}[noitemsep]
    \item Philip, \textit{Zong!}: language fragmented and drowned; Philip does with text what the hydrophone does with water — makes the Atlantic archive audible by fragmenting the dominant narrative \parencite{Philip_2008}.
    \item Neimanis, \textit{Bodies of Water}: trans-corporeality; we are bodies of water, not bodies in it; dissolution as homecoming, not only loss \parencite{Neimanis_2017}.
\end{itemize}

\noindent\textbf{Making Entanglement:} Philip fragments text; this work fragments signal. The formal parallel carries the argument. The ritual of submerging artifacts — salt, sugar, wax, clay — each with its own dissolution rate, sonic character, and symbolic resonance — frames the interaction as compositional chemistry. Neimanis reframes what the hydrophone is listening for: not loss, but transformation. The processing chain has to carry that distinction.

\noindent\textbf{ST Phase:} Remixing. Developing the artifact taxonomy; refining the granular chain to foreground textural evolution rather than discrete event capture.

\medskip

\subsection*{Log Entry 3: Temporal Escape}

\noindent\textbf{Literature/Theory:}
\begin{itemize}[noitemsep]
    \item Phillips / Black Quantum Futurism: Black non-linear temporality; cryosleep as temporal technology — leaving one time, arriving in another, bypassing the colonial timeline \parencite{Phillips_2015}.
    \item Butler, \textit{Xenogenesis / Lilith's Brood}: Lilith wakes from suspension into an alien ship; the Black body simultaneously preserved and violated by what holds it \parencite{Butler_1987}.
    \item Eshun returns: sonic time as a speculative apparatus for escaping the colonial temporal order.
\end{itemize}

\noindent\textbf{Making Entanglement:} NASA sonification samples — electromagnetic waves from space, solar wind rhythms — are mixed with hydrophone recordings from the vessel. The Atlantic and the cosmos occupy the same sonic register. ``Did he fear it?'' demands an answer the historical archive refused to record. The work answers through sonic form: by building a space where earthly dissolution and interstellar passage coexist as the same kind of transit.

\noindent\textbf{ST Phase:} Embodying. Collaborative workshops test whether the sonic environment produces the phenomenology of temporal escape.

\medskip

\subsection*{Log Entry 4: The Instrument Built from This}

\noindent\textbf{Literature/Theory:}
\begin{itemize}[noitemsep]
    \item Barad, \textit{Meeting the Universe Halfway}: the water vessel as apparatus; the sound of dissolution emerges through intra-action of water, artifact, hydrophone, and cultural logic encoded in the processing chain \parencite{Barad_2007}.
    \item Franinović \& Serafin, \textit{Sonic Interaction Design}: interaction design principles applied to ritual submersion as primary interaction modality \parencite{Franinovic_Serafin_2013}.
\end{itemize}

\noindent\textbf{Making Entanglement:} Full technical documentation: hydrophone construction and sourcing, vessel and basin design, artifact taxonomy (dissolution rates, sonic characters, symbolic resonances), Max/MSP spectral analysis and granular synthesis chain, vertical spatialization system (heavier sounds settle toward lower positions, lighter sounds float — organized by fluid dynamics rather than terrestrial spatial norms), conductivity sensor integration (planned). ST cycle applied fully. ST domains applied specifically: FS refuses the separation of Middle Passage from interstellar travel; MM locates water as archive; SW demands a different listening posture from vertical space; IW positions participants as witnesses, and collaborative workshop documentation feeds directly into this domain.

\noindent\textbf{Collaborative Workshop Integration:} Three sessions maximum, three participants maximum, one confirmed. The workshops specifically probe IW: can participants sense the connection between dissolving artifact and dissolving body? Between floating wax and suspended transit? Documentation (audio/video, field notes, reflective journal) is the primary data source for this section.

\noindent\textbf{ST Phase:} Transmitting.

\medskip

\subsection*{Dissolution: What Holds Cannot Be Held}

\noindent\textbf{Making Entanglement:} Synthesis of the movement. The leitmotif here is determined by the chemistry of dissolution rather than a metronome — duration is variable; the material sets the rhythm. The Somax2 corpus draws on water-imagery spirituals and NASA sonifications. The drone shifts timbre: vessel resonance. The movement closes with the question that opens the third: \textit{If water carries the body through time — if dissolution is also transformation — what does the body itself carry? What knowledge lives in the flesh that no algorithm was trained to find?}

\bigskip

% ------------------------------------------------------------
% PART IV: MOVEMENT III — TOUCH
% ------------------------------------------------------------

\section*{Part IV: Movement III --- Touch (\textit{Don't Play with My Hair}, skin/flesh)}
\addcontentsline{toc}{section}{Part IV: Movement III --- Touch}

\noindent\textbf{Tempo:} Human. Natural gesture rhythm, the pace of the body.\\
\noindent\textbf{Form:} Technical manual that keeps interrupting itself. Clinical notation broken by the embodied, the felt, the refused. The form enacts what the piece argues.\\
\noindent\textbf{Material:} CT scan of curly hair (VDB $\to$ point cloud), HairSynth v1 (Houdini + Python + PlugData wavetable) and v2 (vvvv/OpenCV + Max/MSP), gesture recognition, capacitive touch sensors (planned).\\
\noindent\textbf{ST Domain Foregrounded:} Sensory Worlding (SW).\\
\noindent\textbf{Motif Recognition:} Explicit. The developmental cell finally audible as itself.\\
\noindent\textbf{Listener Relationship:} Participated.\\
\noindent\textbf{Development Status:} Adjacent work. Two full technical iterations (HairSynth v1, v2) exist. Prose to be developed from existing scaffolding.

\medskip

\subsection*{Before the Manual Begins: The Refusal and the Instrument}

\noindent\textbf{Literature/Theory:}
\begin{itemize}[noitemsep]
    \item Spillers, ``Mama's Baby, Papa's Maybe'': the distinction between body and flesh; the flesh as the site of racial inscription; the uninvited touch as a claim on the flesh \parencite{Spillers_1987}.
\end{itemize}

\noindent\textbf{Making Entanglement:} The work begins before the instrument. Hands reaching toward my head without permission; the dual experience of my hair as site of violation and site of intimate self-knowledge. ``Don't play with my hair'' is a refusal that also names. The speculative proposition arrives from that lived encounter: what if hair could be a literal instrument? The question makes the HairSynth.

\noindent\textbf{ST Phase:} Unsettling. The desire arises from embodied experience of two failures simultaneously — the social world's treatment of Black hair as spectacle, and the computer vision pipeline's inability to see it.

\medskip

\subsection*{Section 1.0: The Body as Signifying Practice}

\noindent\textbf{Literature/Theory:}
\begin{itemize}[noitemsep]
    \item Mercer, ``Black Hair/Style Politics'': hair as signifying practice; aesthetics and politics as inseparable in the styling of Black hair \parencite{Mercer_1987}.
    \item Spillers returns: the flesh as the site where the HairSynth's reclamation occurs — knowledge and sound generated on the body's own terms.
\end{itemize}

\noindent\textbf{Making Entanglement:} The HairSynth makes Mercer's signifying practice generative in the most literal sense: cultural meaning encoded in hair structure becomes computational input. The sign generates sound. Political context — the CROWN Act, workplace discrimination, school dress codes — establishes that the stakes of this instrument extend beyond the sonic.

\noindent\textbf{ST Phase:} Unsettling. The social and institutional machinery that polices Black hair is the embedded assumption the instrument is built to refuse.

\medskip

\subsection*{Section 2.0: The Pipeline (Or: What the Algorithm Cannot See)}

\noindent\textbf{Literature/Theory:}
\begin{itemize}[noitemsep]
    \item Buolamwini, \textit{Gender Shades}: CV systems fail to detect Black hair because training data encoded whiteness as norm; Type 4 hair is structurally absent from most datasets \parencite{Buolamwini_2018}.
    \item Benjamin, \textit{Race After Technology}: the algorithm's failure is the argument of the algorithm \parencite{Benjamin_2019a}.
\end{itemize}

\noindent\textbf{Making Entanglement:} HairSynth v1: CT scan of a curly fro $\to$ Houdini VDB $\to$ point cloud $\to$ Python script $\to$ WAV files (1024 samples) $\to$ PlugData wavetable synthesizer. HairSynth v2: vvvv + OpenCV $\to$ contour data, keypoints, blobs $\to$ Max/MSP parameter mapping. The technical friction is the research site. The CV system's failure to detect Type 4 hair is the unsettling the work is built upon, not a limitation to route around.

\noindent\textbf{ST Phase:} Remixing. Repurposing CV algorithms to prioritize textural analysis at scales appropriate for hair structure; building the modification on top of the tool's own failure.

\medskip

\subsection*{Interruption I}

\noindent\textbf{Function:} The manual pauses. A passage in a different register — the specific texture of touching my own hair, the sound of clippers at the back of my neck, knowing my hair by feel in the dark. This is the data the pipeline is trying to extract. I have always had it. The pipeline is learning to ask a question my body can already answer. \textit{No citations. Different rhythm. Present tense. Sensory.}

\medskip

\subsection*{Section 3.0: Haptic Epistemology}

\noindent\textbf{Literature/Theory:}
\begin{itemize}[noitemsep]
    \item Marks, \textit{The Skin of the Film}: haptic visuality; the skin as an epistemological organ; touch as a way of knowing that precedes and exceeds vision \parencite{Marks_2000}.
    \item Barad, \textit{Meeting the Universe Halfway}: HairSynth as apparatus; hair data and synthesis system intra-act to produce a sound that neither generates alone \parencite{Barad_2007}.
    \item Grimshaw \& Garner, \textit{Sonic Virtuality}: the drone returns explicitly here — sound was the condition of possibility before any of this was built.
\end{itemize}

\noindent\textbf{Making Entanglement:} The data-sound mappings encode embodied knowledge — what it feels like to touch hair — rather than simple numeric correlation. They emerged through embodied experimentation: adjusting while listening until the output felt right. Marks names why that process is methodology rather than subjectivity. The skin knows before the eye, and the ear confirms what the hand already understood.

\noindent\textbf{ST Phase:} Embodying. Testing whether the synthesis carries the haptic knowledge the body already possesses.

\medskip

\subsection*{Section 4.0: Braiding}

\noindent\textbf{Literature/Theory:}
\begin{itemize}[noitemsep]
    \item Kimmerer, \textit{Braiding Sweetgrass}: braiding as a practice of attending to material with patience, relationship, and intergenerational knowledge \parencite{Kimmerer_2013}. \admargin{Distinct traditions in conversation — handle with care, do not collapse.}
\end{itemize}

\noindent\textbf{Making Entanglement:} The intimate gestures — fingers running through hair, twisting strands, the self-soothing pull — are what the HairSynth is built to receive. Consent becomes a design problem: who touches whose hair, and with what permission? That question does not resolve in this section. The form of the manual, interrupted, enacts the irresolution. The unresolved question is itself transmissible knowledge.

\noindent\textbf{ST Phase:} Embodying $\to$ edges of Transmitting.

\medskip

\subsection*{Section 5.0: The Instrument Built from This}

\noindent\textbf{Literature/Theory:}
\begin{itemize}[noitemsep]
    \item Franinović \& Serafin, \textit{Sonic Interaction Design}: interaction design for a body-driven synthesis system \parencite{Franinovic_Serafin_2013}.
    \item Agre, ``Toward a Critical Technical Practice'': the technical practitioner who turns the discipline's tools on the discipline itself \parencite{Agre_1998}.
\end{itemize}

\noindent\textbf{Making Entanglement:} Full technical documentation of HairSynth v1 and v2: wavetable architecture (PlugData), vvvv contour extraction pipeline, the mapping system (parameters $\leftrightarrow$ synthesis), visual feedback component, planned sensor integration (capacitive touch, proximity sensors). The Somax2 corpus draws on barbershop ambience and hair care oral tradition. Full ST cycle documented. ST domains applied specifically.

\noindent\textbf{ST Phase:} Transmitting.

\medskip

\subsection*{Final Interruption}

\noindent\textbf{Function:} The manual cannot hold its form. The output sounds right — not just technically, but in relation to the experience of having this hair, growing up with it, having it reached for without permission and touching it myself in the mirror. The arc completes: felt (Ground), witnessed (Hold), participated (Touch). The leitmotif arrives in its most recognizable form. The cell is finally audible as itself. \textit{No citations. Short. Let it ring.}

\bigskip

% ------------------------------------------------------------
% PART V: ST ACROSS THE FULL SCORE
% ------------------------------------------------------------

\section*{Part V: How ST Works Across the Full Score}
\addcontentsline{toc}{section}{Part V: ST Across the Full Score}

What follows maps the framework as it operates across all three works simultaneously. You can read any row of the table to see how one domain manifests differently in each material substrate, or read any column to see all four domains at work within a single movement.

\subsection*{The Four Domains Across Three Movements}

\begin{center}
\small
\begin{tabular}{|p{2.8cm}|p{3.5cm}|p{3.5cm}|p{3.5cm}|}
\hline
\textbf{ST Domain} & \textbf{Movement I: Ground} & \textbf{Movement II: Hold} & \textbf{Movement III: Touch} \\
\hline
\textbf{Fugitive Speculation} &
Speculative encounter between soldier and buffalo; refuses colonial positioning of either as simply victim or agent &
Refuses separation of Middle Passage from interstellar travel; ``Did he fear it?'' re-inserts Black subjectivity into suspension narratives &
HairSynth centers Black hair as site of knowledge generation rather than spectacle or problem to manage \\
\hline
\textbf{Material Memory} &
Earth, rock, and soil carry sedimented history; the substrate box is archive and container simultaneously &
Water as archive; dissolving artifacts enact ritual of passage and transformation; the Atlantic and the cosmos as the same space of transit &
Hair carries genetic, environmental, and temporal memory; muscle memory of care practices directly informs gesture mappings \\
\hline
\textbf{Sensory Worlding} &
Tactile interaction as primary input; spatialization informed by migration data rather than concert-hall norms &
Vertical spatialization organized by fluid dynamics; ritual submersion as interaction; hydrophone recordings mixed with NASA sonification &
Haptic-visual-sonic feedback loop; hair's materiality made perceptible across multiple registers simultaneously \\
\hline
\textbf{Intimate Witnessing} &
Does the granular fragmentation feel like rupture? Can participants sense the historical weight encoded in the substrate? &
Collaborative workshops: can witnesses sense dissolution as passage, floating wax as suspended transit? Primary site of IW documentation &
Does the synthesis sound like hair? The final interruption names the moment when the output achieves alignment \\
\hline
\end{tabular}
\end{center}

\subsection*{The URET Cycle Across Three Movements}

Each movement enacts a complete — or substantially complete — cycle of Unsettling $\to$ Remixing $\to$ Embodying $\to$ Transmitting, with different materials, different sites of friction, and different forms of knowledge produced. Taken together, the three cycles make the case that ST is a generalizable methodology, not a single-work proposition.

\begin{itemize}[noitemsep]
    \item \textbf{Ground}: Unsettling at the level of spatial assumption (neutral acoustic space has no memory); Remixing through granular fragmentation and migration-data spatialization; Embodying through tactile interaction testing; Transmitting through documented procedures and code.
    \item \textbf{Hold} (primary): Unsettling at the level of the signal-cleaning default, which erases residue; Remixing through artifact taxonomy and vertical spatialization; Embodying through collaborative witness workshops; Transmitting through documented ritual, workshop findings, and open-source processing chain.
    \item \textbf{Touch}: Unsettling at the level of the CV system's structural failure to see Type 4 hair; Remixing through two full technical iterations (HairSynth v1, v2); Embodying through haptic mapping refined by felt sense; Transmitting through documented pipeline and the unresolved consent design questions.
\end{itemize}

\subsection*{The Through-Line: Transduction and Refraction}

ST operates through two mechanisms that run across all three movements.

\textbf{Transduction}: speculative cultural ideas converted into material form through redesign. In Ground, the speculative encounter between soldier and buffalo becomes granular fragmentation and migration-data spatialization. In Hold, the question about fear and suspension becomes ritual submersion and vertical fluid dynamics. In Touch, the politics of hair and uninvited touch become a haptic synthesis pipeline.

\textbf{Refraction}: altered instruments bending perception toward new sonic and cultural relationships. In Ground, the refracted perception is the ambiguity between stone and flesh. In Hold, it is the co-habitation of dissolution and arrival. In Touch, it is the moment when the synthesis output sounds like what the body already knew.

Together, these two mechanisms constitute the methodological contribution: a replicable, documented process for converting cultural epistemology into sonic and computational form — and for verifying, through practice, that the conversion carried what it was meant to carry.
